\chapter{What Remains Outside the Current Evidence}
\label{ch:outside-current-evidence}

The battery book is stronger when it is explicit about what has \emph{not}
been shown.  This chapter records the outer boundary of the current evidence
package so the thesis does not borrow credibility from work that has not yet
been executed.

\section{Not Claimed}

The current package does not claim:

\begin{enumerate}
  \item universal deployment safety across all cyber-physical domains,
  \item full physical hardware validation on a production battery stack,
  \item adversarial completeness against every possible spoofing strategy,
  \item heterogeneous multi-asset composition guarantees, or
  \item a final certificate-expiration law beyond the conservative blackout
    hold behavior now locked in Chapter~\ref{ch:cert-half-life-blackout}.
\end{enumerate}

\section{Still Partial}

Several thesis-hardening layers exist in meaningful but still bounded form:

\begin{itemize}
  \item Aging and asset preservation are represented through proxy tables and
    calibration design rather than long-horizon field degradation data.
  \item Fleet composition is implemented for the two-battery shared-transformer
    case, not for a broader heterogeneous fleet.
  \item The HIL chapter now has an auditable runtime rehearsal, but not a full
    external-bench package with independent device timing and relay shutdown.
\end{itemize}

\section{Why This Boundary Matters}

The battery thesis is meant to be hard to challenge because it is internally
coherent, not because it quietly widens its scope.  By naming what remains
outside the current evidence, the thesis protects the core battery result from
being diluted by larger claims that the present repo does not yet support.

\section{Summary}

The current evidence package is deep, but it is intentionally not universal.
That is the correct scientific posture for a battery-first ORIUS thesis.
